\section*{6. 酸・アルカリの正体とpH}
\begin{enumerate}
    \item 酸性の水溶液に共通して含まれ、酸性の正体となるイオンを \textbf{化学式（イオン式）} で書きなさい。\\[0.2cm]
    （ \textcolor{red}{$\mathbf{H^+}$} ）
    \item アルカリ性の水溶液に共通して含まれ、アルカリ性の正体となるイオンを \textbf{化学式（イオン式）} で書きなさい。\\[0.2cm]
    （ \textcolor{red}{$\mathbf{OH^-}$} ）
    \item 水溶液の酸性・アルカリ性の強さを表す数値をpHという。下の数直線の(A)～(C)の範囲は、それぞれ何性を示しているか答えなさい。
\end{enumerate}

\begin{center}
\begin{tikzpicture}[scale=0.9]
    \draw[thick, ->] (0,0) -- (15,0);
    \foreach \x in {0, 7, 14} {
        \draw[thick] (\x, 0.2) -- (\x, -0.2) node[below] {\x};
    }
    \draw[<->] (0.2, 0.5) -- (6.8, 0.5) node[midway, above] {(A)};
    \draw[<-] (7, 0.5) -- (7, 1) node[above] {(B)};
    \draw[<->] (7.2, 0.5) -- (14, 0.5) node[midway, above] {(C)};
\end{tikzpicture}
\end{center}
\vspace{0.3cm}
\begin{flushright}
(A)の範囲（ \textcolor{red}{\textbf{酸}} ）性 \hspace{0.5cm} (B)（ \textcolor{red}{\textbf{中}} ）性 \hspace{0.5cm} (C)の範囲（ \textcolor{red}{\textbf{アルカリ}} ）性
\end{flushright}
